\chapter{C\'alculo del \textit{Fitness}}\label{chapter:fitness}

En este capítulo se describe el procedimiento mediante el cual se evalúa la capacidad de un grafo dirigido para modelar un sistema compartimental que represente 
adecuadamente datos observacionales. Cada grafo candidato define de manera única un sistema de ecuaciones diferenciales ordinarias (EDOs), donde los nodos corresponden 
a compartimentos y las aristas representan interacciones entre ellos. A partir de la estructura del grafo y de las trayectorias temporales observadas, se construyen 
las ecuaciones dinámicas correspondientes, ajustando sus coeficientes mediante información extraída de los datos. Esta conversión permite evaluar computacionalmente 
qué tan bien el sistema de EDOs asociado al grafo logra reproducir la dinámica observada, a través de una función de \textit{fitness} que combina precisión en el 
ajuste, plausibilidad estructural y restricciones biológicas.

La sección~\ref{sec:exp_pol} introduce la generación inicial de ecuaciones polinomiales hasta segundo orden para cada compartimento, a partir de una biblioteca de términos 
construida de forma sistemática. En la sección~\ref{sec:term_sig} se detalla el procedimiento de selección de términos significativos, que consiste en ajustar los coeficientes 
por mínimos cuadrados, normalizar sus valores y filtrar aquellos de baja relevancia, manteniendo la coherencia estructural con el grafo. La sección~\ref{sec:cons_sis} 
explica cómo, a partir de los términos seleccionados y las conexiones topológicas del grafo, se construyen las ecuaciones diferenciales completas del sistema dinámico, 
y se describe cómo estas ecuaciones simbólicas se 
convierten en funciones numéricamente evaluables, permitiendo calcular trayectorias simuladas. Finalmente, en la sección~\ref{sec:punt_fit} se presenta la definición completa del 
\textit{fitness}, que incluye el cálculo del error cuadrático medio entre datos y simulaciones, junto con penalizaciones por comportamientos no deseados como 
poblaciones negativas, alta complejidad estructural o sobreconectividad del grafo.

A continuación, se describe el proceso de generación y filtrado de términos polinomiales que sirven como base para construir el sistema dinámico asociado a cada grafo.

\section{Expansión Polinómica del Modelo Dinámico}
\label{sec:exp_pol}

El cálculo del \textit{fitness} comienza con una matriz de datos temporales que describe la evolución de las poblaciones 
en cada compartimento a lo largo del tiempo y un grafo dirigido que representa las interacciones entre dichos compartimentos. 

Para cada compartimento, se generan ecuaciones que contienen todos los posibles términos  polinomiales hasta segundo orden. Esta construcción de una biblioteca de 
funciones candidatas a partir de la cual se seleccionará la estructura del modelo es una práctica establecida en la identificación de sistemas dinámicos no lineales 
\cite{ljung1999, Billings2013}. Los términos considerados son:

\begin{itemize}
    \item Términos lineales del tipo $X_i$, donde $X_i$ es la variable asociada al compartimento $i$.
    \item Términos cuadráticos del tipo $X_i X_j$ para todas las combinaciones posibles de compartimentos $i$ y $j$, incluyendo productos consigo mismos.
\end{itemize}

Por ejemplo, en un sistema con \(X_0, X_1, X_2\) donde cada uno de estos compartimentos representan un nodo del grafo, las ecuaciones asociadas a cada nodo serían:
\begin{align*}
    \frac{dX_0}{dt} &= a_{00} X_0 + a_{01} X_1 + a_{02} X_2 + a_{03} X_0^2 + a_{04} X_0 X_1 + a_{05} X_0 X_2 + a_{06} X_1^2 + a_{07} X_1 X_2 + a_{08} X_2^2, \\
    \frac{dX_1}{dt} &= a_{10} X_0 + a_{11} X_1 + a_{12} X_2 + a_{13} X_0^2 + a_{14} X_0 X_1 + a_{15} X_0 X_2 + a_{16} X_1^2 + a_{17} X_1 X_2 + a_{18} X_2^2, \\
    \frac{dX_2}{dt} &= a_{20} X_0 + a_{21} X_1 + a_{22} X_2 + a_{23} X_0^2 + a_{24} X_0 X_1 + a_{25} X_0 X_2 + a_{26} X_1^2 + a_{27} X_1 X_2 + a_{28} X_2^2.
\end{align*}

La inclusión sistemática de estos términos se inspira en técnicas clásicas de modelado de sistemas dinámicos, donde los términos lineales permiten representar 
efectos proporcionales como tasas constantes de crecimiento o pérdida \cite{strogatz2018nonlinear}, mientras que los términos cuadráticos capturan interacciones 
no lineales entre compartimentos, tales como competencia \cite{murray2002mathematical}, cooperación o transmisión \cite{anderson1992infectious}. 

Este conjunto base de funciones puede entenderse como una expansión polinómica de orden dos, similar a una serie de Taylor truncada \cite{ArfkenWeberHarris2012}, 
diseñada para capturar la dinámica local del sistema mediante combinaciones algebraicas con complejidad acotada. Al incluir desde el inicio todos los términos 
polinomiales de primer y segundo orden que puedan influir en la evolución del sistema, se evita descartar prematuramente posibles interacciones relevantes.

Este paso proporciona una base estructurada sobre la cual se realiza, en la próxima sección, la selección de los términos significativos que son aquellos que 
resultan relevantes para describir la dinámica observada.

\section{Selección de Términos Significativos}
\label{sec:term_sig}

Una vez generadas las ecuaciones expandidas para cada compartimento, el siguiente paso consiste en ajustar sus coeficientes y seleccionar los términos significativos. 
La relevancia de un término se determina en función de la magnitud y el signo de su coeficiente tras el ajuste, así como de su coherencia con la topología del grafo. 
Este proceso se desarrolla en tres etapas: ajuste de coeficientes, eliminación de términos no significativos, e identificación de términos compartidos entre 
compartimentos conectados.

Para cada ecuación, se estiman las derivadas numéricas de los datos originales utilizando diferencias finitas hacia adelante, lo que permite aproximar las 
tasas de cambio temporales. Por ejemplo, si se dispone de datos de población para el compartimento \(X_1\) 
en los tiempos \(t_0\) y \(t_1\), la derivada en \(t_0\) se puede calcular como:
\[
\dot{X}_1(t_0) \approx \frac{X_1(t_1) - X_1(t_0)}{\Delta t}.
\]
Para el último punto de la serie temporal, se utiliza una fórmula de diferencia finita hacia atrás:
\[
\dot{X}_1(t_{n-1}) \approx \frac{X_1(t_{n-1}) - X_1(t_{n-2})}{\Delta t}.
\]

A continuación, se construye una matriz que contiene la evaluación de cada término algebraico en cada instante de tiempo. Esta matriz permite reformular el 
problema como un sistema lineal, donde los coeficientes de la ecuación diferencial son las incógnitas, y la derivada numérica actúa como vector objetivo. 
Este enfoque es común en la identificación de sistemas dinámicos \cite{ljung1999, Billings2013}.

Por ejemplo, si una ecuación incluye los términos \(X_0\), \(X_1\) y \(X_0 X_1\), y se tienen observaciones en 4 puntos de tiempo, entonces la matriz de diseño 
tendrá 4 filas y 3 columnas, donde cada entrada corresponde al valor del término evaluado en ese instante. El sistema se resuelve mediante mínimos cuadrados para 
encontrar los coeficientes que minimizan la diferencia entre las derivadas observadas y las predichas.

Los coeficientes obtenidos se normalizan al intervalo \([0, 1]\), conservando el signo original. Esta normalización facilita la comparación entre términos 
y permite interpretar la magnitud relativa de cada interacción \cite{Bishop2006, Hastie2009}.  Preservar los signos es importante, ya que indica la dirección del 
flujo: coeficientes positivos representan aportes a un compartimento, mientras que valores negativos corresponden a salidas.

Tras el ajuste, se eliminan todos los términos cuyo coeficiente sea inferior a un umbral predefinido. Este umbral actúa como un hiperparámetro que regula el equilibrio 
entre la complejidad del modelo y su capacidad de ajuste. Por ejemplo, si el umbral es 0.1, el término $0.05 X_2^2$ sería eliminado por considerarse no significativo. 
Un umbral bajo puede conducir a modelos con demasiados términos irrelevantes y, por ende, propensos al sobreajuste \cite{Hastie2009}. Por el contrario, un umbral 
excesivamente alto podría descartar interacciones relevantes.  Esta técnica de eliminación de términos basada en la magnitud de sus coeficientes es una forma de 
selección de variables o reducción de la complejidad del modelo \cite{Billings2013, Hastie2009}.

La eliminación de t\'erminos no significativos responde al principio de parsimonia en modelado matemático \cite{Burnham2002, Box1976}, priorizando modelos más simples 
y comprensibles sin sacrificar precisión. El resultado es un conjunto reducido de términos con impacto significativo en la dinámica poblacional observada.

Después de eliminar los términos no relevantes, se comparan las ecuaciones de los nodos que están conectados entre sí en el grafo. Se identifican los términos que 
aparecen simultáneamente en ambas ecuaciones; por ejemplo, si los compartimentos \(X_1\) y \(X_2\) están conectados, sus ecuaciones diferenciales respectivas luego 
del ajuste y selección de términos significativos quedan como:
\begin{align*}
\frac{dX_1}{dt} &= 0.5 X_1 + 0.2 X_1 X_2 + 0.1 X_2^2, \\
\frac{dX_2}{dt} &= -0.4 X_2 + 0.3 X_1 X_2 + 0.2 X_1^2.
\end{align*}
El término \(X_1 X_2\) aparece en ambas ecuaciones, indicando una interacción directa entre los compartimentos. Por tanto, este término 
se conserva como representativo de la influencia mutua modelada por la arista entre \(X_1\) y \(X_2\).

En caso de que no existan términos comunes entre dos nodos conectados, se introduce al menos un término que combine sus variables, como \(X_i X_j\). Esto garantiza 
que toda conexión topológica en el grafo tenga una contraparte algebraica en el sistema de ecuaciones. Así se mantiene la coherencia entre la estructura del grafo y la 
formulación matemática del modelo. Esta adición no es arbitraria: el término \(X_i X_j\) tiene una interpretación natural como interacción entre las poblaciones de 
ambos compartimentos. Por ejemplo, \(X_0\) y \(X_3\) están conectados en el grafo, pero después del filtrado de términos significativos no comparten ningún término y 
sus respectivas ecuaciones quedan como las siguientes: 
\begin{align*}
\frac{dX_0}{dt} &= 0.6 X_0 - 0.2 X_1^2, \\
\frac{dX_3}{dt} &= -0.5 X_3 + 0.1 X_2 X_3.
\end{align*}
Entonces, para asegurar la representación dinámica de esa conexión, se incorpora el término \(X_0 X_3\) como representativo en ambas ecuaciones, que refleja la 
interacción entre \(X_0\) y \(X_3\), en concordancia con la arista existente entre ellos en el grafo. De esta manera, se evita que la estructura topológica pierda 
validez dinámica en el modelo.

Este mecanismo garantiza que las conexiones entre compartimentos deben estar respaldadas por interacciones en las ecuaciones que reflejen adecuadamente los posibles 
flujos o influencias entre ellos.

El siguiente paso consiste en construir formalmente el sistema dinámico completo. Esta reconstrucción se realiza a partir de los elementos obtenidos hasta ahora: 
los términos significativos entre las ecuaciones, las conexiones del grafo y los coeficientes ajustados. En la próxima sección se describe en detalle cómo se generan 
las ecuaciones diferenciales a partir de estos componentes.

\section{Construcción del Sistema Dinámico Basado en el Grafo}
\label{sec:cons_sis}

Despu\'es de seleccionados los términos significativos, se procede a construir el sistema dinámico definitivo, en el cual las ecuaciones de cada compartimento se ajustan 
según la estructura topológica del grafo propuesto. Este proceso consta de tres etapas: la reconstrucción estructurada de las ecuaciones y el reajuste de coeficientes 
en el nuevo modelo y la conversión del sistema en funciones computables que permitan calcular las derivadas en cada instante de tiempo como parte de una integración 
diferencial mediante métodos numéricos.

Para cada compartimento, se genera una ecuación que incluye únicamente aquellos términos derivados de interacciones con otros nodos a los que está conectado 
directamente mediante una arista. La dirección de la arista influye directamente en la construcción de las ecuaciones: si el flujo va desde otro nodo hacia el 
nodo actual, los términos asociados tienen coeficientes positivos, indicando una contribución entrante; en cambio, si la arista sale del nodo, los coeficientes 
correspondientes son negativos, reflejando una pérdida o salida desde ese compartimento. Si el nodo \(X_1\) está conectado con \(X_0\) y \(X_2\) a través de las 
aristas \(X_0 \rightarrow X_1\) y \(X_1 \rightarrow X_2\), entonces su ecuación podrá contener términos como \(+a X_0\), \(+b X_0 X_1\) y \(-c X_1\), \(-d X_1 X_2\). 
Así, la ecuación de \(X_1\) podría tomar la forma:
\[
\frac{dX_1}{dt} = a X_0 + b X_0 X_1 - c X_1 - d X_1 X_2,
\]
donde los coeficientes \(a\), \(b\), \(c\), y \(d\) son positivos, y el signo con el que aparecen en la ecuación refleja si la interacción representa una ganancia 
o una pérdida para el nodo \(X_1\).

El modelo resultante mantiene así la consistencia estructural del grafo, en el sentido de que ningún nodo es influido por otros si no existe una conexión directa, 
lo que proporciona interpretabilidad y transparencia a las interacciones modeladas.

Tras la reconstrucci\'on de las ecuaciones, se realiza un nuevo ajuste de coeficientes mediante m\'etodos de m\'inimos cuadrados. Esta regresi\'on se aplica solo 
sobre el subconjunto de t\'erminos seleccionados y permitidos por la topolog\'ia, a fin de refinar los par\'ametros y mejorar el ajuste a las derivadas observadas.
Por ejemplo, si la ecuaci\'on de $X_1$ se reduce a $\dot{X}_1 = aX_0 - bX_1X_2$, los valores de $a$ y $b$ se recalculan para minimizar la diferencia entre la 
derivada observada y la predicha. Este enfoque de selecci\'on seguida de reajuste se asemeja a t\'ecnicas como \textit{LASSO} \cite{Tibshirani1996} o regresi\'on 
\textit{stepwise} \cite{Efroymson1960}, pero adaptado al contexto estructural impuesto por el grafo. As\'i se equilibra la interpretabilidad del modelo con su capacidad 
predictiva.

Un ejemplo de un grafo con tres nodos \(X_0\), \(X_1\), y \(X_2\), conectados mediante las aristas \(X_0 \rightarrow X_1\), \(X_1 \rightarrow X_2\), y 
\(X_0 \rightarrow X_2\). Tras el proceso de selección, reconstrucción y reajuste de coeficientes descrito, las ecuaciones del sistema dinámico podrían 
tomar la siguiente forma:
\begin{align*}
\frac{dX_0}{dt} &= -0.2 X_0 - 0.1 X_0 X_1 - 0.4 X_2, \\
\frac{dX_1}{dt} &= 0.2 X_0 + 0.1 X_0 X_1 - 0.6 X_1 - 0.3 X_1 X_2, \\
\frac{dX_2}{dt} &= 0.4 X_2 + 0.6 X_1 + 0.3 X_1 X_2.
\end{align*}
En este ejemplo, los signos reflejan la dirección del flujo: \(X_0\) pierde población hacia \(X_1\) y \(X_2\), lo que se refleja en términos negativos en su 
ecuación. A su vez, \(X_1\) recibe flujo desde \(X_0\) y lo envía hacia \(X_2\), resultando en una combinación de términos positivos y negativos. Finalmente, 
\(X_2\) actúa como sumidero en este grafo, acumulando los flujos entrantes desde \(X_0\) y \(X_1\).

Una vez definidos los términos y reajustados sus coeficientes, las ecuaciones se convierten de su forma simbólica a funciones evaluables numéricamente. Esta conversión 
es necesaria porque permite calcular el valor de las derivadas en cada instante de tiempo, lo cual es fundamental para poder integrar el sistema de ecuaciones y 
comparar los resultados con los datos observados. Para la manipulación algebraica se utiliza la biblioteca \texttt{SymPy} \cite{SymPy2017}, y su función 
\texttt{lambdify} permite transformar las expresiones simbólicas en funciones eficientes que pueden ser evaluadas directamente por algoritmos numéricos.

El resultado es un sistema de ecuaciones diferenciales ordinarias que puede ser evaluado computacionalmente. Esta capacidad es clave para calcular el error cuadrático 
medio (\textit{mean squared error}, MSE), el cual compara las soluciones del sistema con los datos reales. Dicho error se usa para asignar un puntaje de \textit{fitness} 
al sistema de ecuaciones construido, cuantificando qué tan bien el modelo representado por el grafo de entrada logra reproducir la dinámica observada. A continuación, 
se describe detalladamente este proceso de puntaje del \textit{fitness}.

\section{Puntaje del \textit{Fitness}}
\label{sec:punt_fit}

Ya construidas las ecuaciones diferenciales para cada compartimento, se evalúa su capacidad para reproducir los datos observados. Este proceso se realiza mediante 
integración numérica del sistema y posterior comparación entre las trayectorias simuladas y los datos originales. El resultado se expresa en un puntaje de 
\textit{fitness} que guía el proceso evolutivo.

El sistema de ecuaciones diferenciales se resuelve numéricamente utilizando la función \texttt{solve\_ivp} de la biblioteca \texttt{SciPy} \cite{Virtanen2020SciPy}. El método inicial empleado es 
\texttt{RK45}.

Una vez obtenidas las trayectorias simuladas, se calcula el error cuadrático medio (MSE) entre estas trayectorias y los datos originales. El MSE es especialmente 
adecuado como métrica de comparación porque penaliza de manera más severa las grandes desviaciones, lo que lo convierte en una función objetivo sensible al ajuste
y es ampliamente utilizado en problemas de regresión y ajuste de modelos \cite{Hastie2009, Bishop2006}.

A este error base se le añaden penalizaciones que incorporan restricciones biológicas, estructurales y computacionales, con el fin de guiar el proceso hacia modelos 
más plausibles. Las penalizaciones son las siguientes:

\begin{itemize}
    \item \textbf{No conservación de la población total:} Se penalizan soluciones en las que la población total varía, solo en caso de que el par\'ametro de 
    esta penalizaci\'on sea activado, bajo la hipótesis de un sistema cerrado sin nacimientos ni migraciones. La penalización se calcula como la varianza de la 
    masa total normalizada por el valor inicial:
    \[
    \text{penalización}_\text{masa} = \frac{\operatorname{Var}\left( \sum_i X_i(t) \right)}{\sum_i X_i(0)} \cdot \lambda_{\text{masa}}, \quad \lambda_{\text{masa}} = 10.
    \]

    \item \textbf{Presencia de valores negativos:} Se penaliza la aparición de poblaciones negativas, las cuales carecen de sentido físico. La penalización corresponde a la suma de los cuadrados de los valores negativos:
    \[
    \text{penalización}_\text{negativos} = \sum_{i,t} \min\left(X_i(t), 0\right)^2 \cdot \lambda_{\text{negativos}}, \quad \lambda_{\text{negativos}} = 5.
    \]

    \item \textbf{Complejidad estructural:} Se penaliza el número de aristas del grafo para favorecer modelos más simples. La penalización depende de la proporción cuadrática de aristas respecto al máximo posible:
    \[
    \text{penalización}_\text{estructura} = \text{MSE} \cdot \left(1 + \alpha \left( \frac{n_\text{aristas}}{n(n-1)} \right)^2 \right), \quad \alpha = 0.5.
    \]

    \item \textbf{Conectividad total del grafo:} Si el grafo tiene el número máximo de aristas dirigidas posibles, se aplica una penalización multiplicativa:
    \[
    \text{penalización}_\text{total} = \text{MSE} \cdot \beta, \quad \text{si } n_\text{aristas} = n(n-1), \quad \beta = 1.5.
    \]
\end{itemize}

El puntaje de \textit{fitness} final se define como la suma del error cuadrático medio y las penalizaciones aplicadas. Este valor refleja cuán bien el modelo describe 
los datos, considerando tanto precisión como simplicidad estructural. Un menor puntaje indica mejor desempeño.

En este capítulo se presentó el procedimiento detallado para evaluar la calidad de un modelo compartimental representado por un grafo dirigido. Se describió cómo, a 
partir de datos, se construyen ecuaciones diferenciales asociadas a cada compartimento mediante una combinación sistemática de términos polinomiales, 
los cuales se filtran en función de su relevancia. Luego, se explicó cómo se reconstruyen las ecuaciones conforme a la topología del grafo y 
cómo se ajustan los coeficientes de forma coherente con las direcciones del flujo. Finalmente, se definió una función de \textit{fitness} que cuantifica el grado de 
ajuste del sistema dinámico a los datos observados, incorporando penalizaciones que reflejan restricciones estructurales, biológicas y computacionales. Esta función 
actúa como el mecanismo de evaluación central para comparar distintas configuraciones de modelos compartimentales.

Contando ya con un procedimiento que permite asignar un puntaje de \textit{fitness} a cualquier sistema de ecuaciones diferenciales generado a partir de un grafo 
compartimental, el siguiente desafío consiste en explorar de manera eficiente el espacio de todos los posibles grafos que podrían representar la dinámica subyacente. 
Este problema de búsqueda estructural es abordado en el próximo capítulo mediante un algoritmo genético, una técnica evolutiva inspirada en la selección natural. A 
través de operadores como mutación, cruzamiento y selección, el algoritmo permite generar, modificar y evaluar poblaciones de grafos candidatos. El objetivo es 
identificar automáticamente estructuras de modelos que modele los datos observados.